\section{Integrated moments of orthonormal projections}
\label{p:projection}

This section answers the first question: how does a bound on the mass of
every Euclidean ball control the density after a random orthonormal
projection, when the label is integrated over its whole space? The proof
has two geometric parts. A frame calculation expresses the cost of
projecting several points through their affine Gram determinant. A covering
argument then integrates that determinant using the ball-mass bound.
We retain the dimensional factor in the density norm because it cancels
against the projected support volume in the regression application.

For integers \(1\le k<d\), let \(\mathcal V_{d,k}\) be the space of
\(k\)-by-\(d\) matrices \(P\) satisfying \(PP^{\mathsf T}=I_k\).
Write \(\vartheta_{d,k}\) for the law of the first \(k\) rows of a uniform
orthogonal matrix. One direct construction fixes the probability
normalization. For a square standard Gaussian matrix \(G\), which is
invertible almost surely, put
\(O=G(G^{\mathsf T}G)^{-1/2}\). Then \(O\) is orthogonal. For fixed
orthogonal \(U,V\), the same formula applied to \(UGV\) gives \(UOV\);
the Gaussian law is unchanged by both multiplications. Thus the law of
\(O\) is a probability invariant on both sides, the Haar probability
used here. Its columns can equivalently be sampled successively uniformly
in the orthogonal complements of the preceding columns. To verify the
equivalence, average either invariant frame law under an independent \(O\);
transitivity makes the averaged law independent of the initial frame.
As in Section~\ref{p:geometry-model}, \(v_k\) is the unit-ball volume.

\subsection{The projected frame and its determinant}

We first record the classical truncated Haar density, together with the
affine-height form used below. Jiang~\cite[Lemma~2.5]{Jiang2006} records
this formula and refers to earlier formulations by Diaconis, Eaton, and
Lauritzen and by Eaton. We include its derivation to retain the
dimensional normalization.

\begin{lemma}[Projected frame density]\label{p:frame-density}
Let \(k,m\ge1\) and \(k+m+1\le d\). Fix a \(d\)-by-\(m\) matrix \(E\)
with orthonormal columns, and draw \(P\sim\vartheta_{d,k}\). The
\(k\)-by-\(m\) matrix \(PE\) has Lebesgue density
\begin{equation}\label{p:frame-density-exact}
 \left(\prod_{j=0}^{m-1}
 \frac{\Gamma((d-j)/2)}
 {\pi^{k/2}\Gamma((d-j-k)/2)}\right)
 \det(I_k-ZZ^{\mathsf T})^{(d-k-m-1)/2}
 \1_{\{I_k-ZZ^{\mathsf T}>0\}}.
\end{equation}
Here \(>0\) means positive definite. In particular, this density is
bounded by \((C\sqrt d)^{km}\).

If \(x_1,\ldots,x_{m+1}\in\R^d\) are affinely independent and
\[
 \delta_i=\dist\bigl(x_i,\aff(x_1,\ldots,x_{i-1})\bigr),
 \qquad 2\le i\le m+1,
\]
then the joint Lebesgue density of the columns
\(P(x_i-x_1)\), \(2\le i\le m+1\), is bounded by
\begin{equation}\label{p:frame-difference-density}
 (C\sqrt d)^{km}\prod_{i=2}^{m+1}\delta_i^{-k}.
\end{equation}
\end{lemma}

\begin{proof}
Right rotational invariance allows us to replace the fixed columns of
\(E\) by an independent uniform orthonormal frame. Conditional on \(P\),
its projection has the same law as the first \(k\) coordinates of that
frame. Sample the full columns successively and let \(Z_j\) be the matrix
of the first \(k\) coordinates of its first \(j\) columns. Set
\[
 S_j=I_k-Z_jZ_j^{\mathsf T},\qquad S_0=I_k.
\]
Given the preceding full columns, the next column is uniform on the unit
sphere in their \(d-j\)-dimensional orthogonal complement. The coordinate
map from that complement to the first \(k\) coordinates has row Gram
matrix \(S_j\). When \(S_j>0\), multiplication by \(S_j^{-1/2}\) makes
its rows orthonormal. Lemma~\ref{p:coordinate-density} and a linear change
of variables therefore give the next projected column \(z\) conditional
density
\[
 c(d-j,k)(\det S_j)^{-1/2}
 (1-z^{\mathsf T}S_j^{-1}z)^{(d-j-k-2)/2}
\]
on \(z^{\mathsf T}S_j^{-1}z<1\). This expression depends only on \(Z_j\),
so it is also the conditional density given those projected columns.
The strict inequality implies \(S_{j+1}=S_j-zz^{\mathsf T}>0\)
almost surely. The determinant identity
\[
 \det S_{j+1}=\det S_j(1-z^{\mathsf T}S_j^{-1}z)
\]
shows that multiplication of the \(m\) conditional densities cancels
every intermediate determinant. The remaining factor is
\(\det S_m^{(d-k-m-1)/2}\), proving \eqref{p:frame-density-exact}.
Its exponent is nonnegative and \(0<\det S_m\le1\).

The normalizing constants have the stated scale. For a unit-scale gamma
variable \(G\) of shape \(a=(d-j-k)/2>0\), Cauchy--Schwarz and the gamma
recurrence give
\[
 \frac{\Gamma((d-j)/2)}{\Gamma((d-j-k)/2)}
 =\E G^{k/2}
 \le(\E G^k)^{1/2}
 =\left(\prod_{i=0}^{k-1}(a+i)\right)^{1/2}
 \le d^{k/2}.
\]
Together with the factor \(\pi^{-k/2}\), this bounds each \(c(d-j,k)\)
by \((C\sqrt d)^k\).

For the last assertion, Gram--Schmidt writes
\([x_2-x_1,\ldots,x_{m+1}-x_1]=ER\), where \(E\) has orthonormal
columns and \(R\) is triangular with
\(\abs{\det R}=\prod_{i=2}^{m+1}\delta_i\). Right multiplication by \(R\)
on \(k\)-by-\(m\) matrices has Jacobian \(\abs{\det R}^k\), one factor
for each row. Applying this change of variables to the preceding density
proves \eqref{p:frame-difference-density}.
\end{proof}

The product of the affine heights in \eqref{p:frame-difference-density}
is \(m!\) times the \(m\)-dimensional volume of the simplex spanned by
the points. Related simplex-volume factors occur in Drury's affine-plane
Jacobian~\cite[Lemma~1, pp.~497--498]{Drury1984}. Here the exponent and
normalization come from the displayed \(k\)-row change of variables.

\subsection{From local mass to an exact density norm}

The local growth hypothesis below is of Frostman type, and
inverse-distance averaging is a classical projection method; compare
Mattila~\cite[Lemma~5.1, p.~236, and Theorem~5.4, p.~237]{Mattila1975}.
We prove the quantitative affine-tube and growing-order moment bounds
needed here.

\begin{lemma}[Affine inverse distances]\label{p:affine-inverse}
Let \(0<\beta\le d\), \(K\ge0\), and \(R>0\). Let \(\nu\) be a finite
positive Borel measure supported in \(B(z,R)\subset\R^d\) such that
\begin{equation}\label{p:all-ball-mass}
 \nu(B(b,t))\le Kt^\beta
 \qquad(b\in\R^d,\ t>0).
\end{equation}
If \(F\) is an affine \(j\)-plane, \(k\ge1\) is an integer, and
\(\beta-j-k\ge1\), then \(\nu(F)=0\) and
\begin{equation}\label{p:affine-inverse-bound}
 \int\dist(x,F)^{-k}\,d\nu(x)\le C^d K R^{\beta-k},
\end{equation}
provided \(j,k\le d\).
\end{lemma}

\begin{proof}
For \(0<t\le R\), the projection of \(B(z,R)\) onto \(F\) is contained
in a \(j\)-ball of radius \(R\). A maximal \(t\)-separated subset is a
\(t\)-net with at most \((3R/t)^j\) points, by packing disjoint
\(j\)-balls of radius \(t/2\). For \(j=0\), one point suffices.
The corresponding ambient balls of radius \(2t\) cover the points of
\(B(z,R)\) within distance \(t\) of \(F\). Thus
\begin{equation}\label{p:affine-tube}
 \nu\{\dist(x,F)\le t\}\le 3^j2^\beta K R^j t^{\beta-j}.
\end{equation}
Because \(\beta-j>0\), continuity from above as \(t\downarrow0\)
gives \(\nu(F)=0\). Distances greater than \(R\) contribute at most
\(R^{-k}\nu(B(z,R))\le KR^{\beta-k}\). On the shell
\(2^{-h-1}R<\dist(x,F)\le2^{-h}R\), \eqref{p:affine-tube} gives a
contribution at most
\[
 3^j2^{\beta+k}KR^{\beta-k}\,2^{-h(\beta-j-k)}.
\]
The sum over \(h\ge0\) is at most twice its first term. Since
\(\beta,j,k\le d\), the resulting factor is bounded by \(C^d\).
\end{proof}

\begin{theorem}[Integrated orthonormal-projection moment]
\label{p:frame-moment}
Let \(d,k,q\) be integers and \(\beta\) a real number with
\(k\ge1\), \(q\ge2\), and
\begin{equation}\label{p:frame-margins}
 k+q\le d,\qquad 0<\beta\le d,\qquad
 \beta-k-(q-2)\ge1.
\end{equation}
Let \(\nu\) be a finite positive Borel measure supported in
\(B(z,R)\), \(R>0\), satisfying \eqref{p:all-ball-mass} for some
\(K\ge0\). Assume that \(P_\#\nu\) is absolutely continuous with
respect to Lebesgue measure for \(\vartheta_{d,k}\)-almost every \(P\).
There is a jointly measurable version \(g(P,u)\) of its density such that
\begin{equation}\label{p:frame-moment-bound}
 \left(\int_{\mathcal V_{d,k}}\int_{\R^k}g(P,u)^q\,du\,
 d\vartheta_{d,k}(P)\right)^{1/q}
 \le C^d(C\sqrt d)^{k(1-1/q)}
 K R^{\beta-k(1-1/q)}.
\end{equation}
For every \(P\), this version vanishes when \(\norm{u-Pz}>R\).
\end{theorem}

\begin{proof}
When \(K=0\), the measure is zero and the assertion is immediate.
For \(K>0\), let \(v_k\) be the unit-ball volume and define
\[
 g_\delta(P,u)=\frac{\nu\{x:\norm{Px-u}\le\delta\}}{v_k\delta^k}.
\]
These are jointly measurable nonnegative functions. Expanding their
\(q\)th powers and using Tonelli's theorem expresses the integrated
moment through \(q\) points \(x_1,\ldots,x_q\) drawn from \(\nu\).
The intersection of the \(q\) label balls of radius \(\delta\) has
volume at most \(v_k\delta^k\); it is empty unless
\(\norm{P(x_i-x_1)}\le2\delta\) for \(i\ge2\).

Put
\(\delta_i=\dist(x_i,\aff(x_1,\ldots,x_{i-1}))\).
Lemma~\ref{p:affine-inverse} shows that every affine plane encountered
here has zero \(\nu\)-mass: its dimension is at most \(q-2\), so the
last condition in \eqref{p:frame-margins} applies. Consequently the tuple
is affinely independent for \(\nu^{\otimes q}\)-almost every choice.
For such a tuple, Lemma~\ref{p:frame-density} applies with \(m=q-1\);
the condition \(k+q\le d\) is exactly its determinant condition. The
probability of all \(q-1\) inequalities is bounded by
\[
 (C\sqrt d)^{k(q-1)}
 \prod_{i=2}^q\delta_i^{-k}
 \bigl[v_k(2\delta)^k\bigr]^{q-1}.
\]
The label-ball volumes cancel those in the definition of \(g_\delta\).
Absorbing \(2^{k(q-1)}\) in \(C\), we obtain
\begin{equation}\label{p:frame-before-integration}
 \int\!\!\int g_\delta^q\,du\,d\vartheta_{d,k}
 \le (C\sqrt d)^{k(q-1)}
 \int\prod_{i=2}^q\delta_i^{-k}\,d\nu^{\otimes q}.
\end{equation}
The degenerate tuples form a null set and can be assigned an infinite
upper bound in this nonnegative integral.

Integrate first in \(x_q\), then in \(x_{q-1},\ldots,x_2\).
At each step Lemma~\ref{p:affine-inverse} bounds the current factor by
\(C^d K R^{\beta-k}\). The remaining first-point mass is at most
\(KR^\beta\). Hence the last integral in
\eqref{p:frame-before-integration} is at most
\[
 C^{d(q-1)}K^qR^{q\beta-k(q-1)}.
\]
This proves the \(q\)th power of \eqref{p:frame-moment-bound} uniformly
in \(\delta\).

Define \(g(P,u)=\liminf_{j\to\infty}g_{1/j}(P,u)\). The assumed absolute
continuity and Lebesgue differentiation
\cite[Theorem~1.6.19, pp.~146--147]{TaoMeasure2011} imply that, for almost every
\(P\), this is a density of \(P_\#\nu\) for almost every \(u\).
Fatou's lemma on the product space
\cite[Corollary~1.4.47, p.~110]{TaoMeasure2011} gives
\eqref{p:frame-moment-bound}. If \(\norm{u-Pz}>R\), sufficiently small
label balls miss \(P(B(z,R))\), since \(P\) is a contraction; the
displayed version is then zero. This limiting calculation concerns the
exact projected measure. It introduces no noise into an observation.
\end{proof}

We now spell out the two parameter choices that will be used. Their
normalizations differ, even though their frame calculation is shared.

\begin{corollary}[All-radii projection of spherical local mass]
\label{p:frame-n30}
For sufficiently large \(d\), put
\[
 n=d-1,\qquad k=q=\lfloor d/4\rfloor,\qquad \beta=n-q.
\]
Let \(0\le f\le1\) be Borel on \(S^{d-1}\), and suppose for \(K\ge0\) that
\[
 \int_{B(b,t)}f\,d\sigma\le Kt^\beta
 \qquad(b\in\R^d,\ t>0).
\]
For \(z\in\R^d\), \(R>0\), and \(E=B(z,R)\), the projection of \(\1_E f\sigma\) by
\(P\sim\vartheta_{d,k}\) has a jointly measurable Lebesgue density
\(g(P,u)\) satisfying
\begin{equation}\label{p:frame-n30-bound}
 \norm g_{L^q(\vartheta_{d,k}\,du)}
 \le C^d(C\sqrt d)^{k(1-1/q)}
 K R^{\beta-k(1-1/q)}.
\end{equation}
\end{corollary}

\begin{proof}
Restriction preserves the stated all-ball bound. For every \(P\), the
projection of \(\sigma\) is absolutely continuous by
Lemma~\ref{p:coordinate-density}, so the same holds for the restricted
weighted measure. The exact margins are
\[
 k+q=2\lfloor d/4\rfloor\le d,\qquad
 \beta-k-(q-2)=d+1-3\lfloor d/4\rfloor\ge d/4+1.
\]
Thus Theorem~\ref{p:frame-moment} applies. This also records the
nonnegative frame-determinant exponent
\((d-k-q)/2=(d-2\lfloor d/4\rfloor)/2\).
\end{proof}

\begin{theorem}[Normalized moment for unit-ball measures]
\label{p:frame-positive}
For sufficiently large \(d\), put
\[
 k=p=\lfloor d/16\rfloor,\qquad \lambda=(d-1)/2.
\]
Let \(X\) be a \(k\)-by-\(d\) standard Gaussian matrix and define
\[
 H=(XX^{\mathsf T})^{1/2},\qquad P=H^{-1}X
\]
on its probability-one full-row-rank event, and set \(P\) to a fixed frame
on the null complement. Let \(m_k\) be uniform
probability on the unit ball of \(\R^k\). Suppose \(\nu\) is a finite
positive Borel measure fixed independently of \(X\), of mass at most one, supported in the unit ball
of \(\R^d\), such that for a fixed \(D>0\),
\[
 \nu(B(b,t))\le D^d t^\lambda
 \qquad(b\in\R^d,\ t>0).
\]
Assume \(P_\#\nu\) is absolutely continuous for almost every \(X\).
There is a jointly measurable nonnegative \(g_X(z)\) such that, for almost
every \(X\), \(P_\#\nu(dz)=g_X(z)\,dz/v_k\). This version vanishes when
\(\norm z>1\) and satisfies
\begin{equation}\label{p:frame-positive-bound}
 \E_X\int g_X(z)^p\,dm_k(z)\le \exp(C_Ddp).
\end{equation}
\end{theorem}

\begin{proof}
Right rotational invariance of \(X\) makes \(P\) have law
\(\vartheta_{d,k}\); the independence assertion about \(P,H\) used later
is proved below. The exact determinant and inverse-distance margins are
\[
 k+p=2\lfloor d/16\rfloor\le d,\qquad
 \lambda-k-(p-2)
 =\frac{d-1}{2}-2\lfloor d/16\rfloor+2
 \ge\frac{3d}{8}+\frac32.
\]
Apply Theorem~\ref{p:frame-moment} with \(q=p\), \(\beta=\lambda\),
\(K=D^d\), and \(R=1\). If \(\ell(P,z)\) is its Lebesgue density,
then \(g_X(z)=v_k\ell(P,z)\) is the required normalized density. Moreover
\[
 \left(\E_X\int g_X^p\,dm_k\right)^{1/p}
 =v_k^{1-1/p}\norm\ell_{L^p(\vartheta_{d,k}\,dz)}.
\]
Use the unit-ball volume estimate \eqref{p:unit-ball-volume}.
Since \(k\ge d/32\) for sufficiently large \(d\),
\[
 v_k^{1-1/p}(C\sqrt d)^{k(1-1/p)}\le C^d.
\]
Theorem~\ref{p:frame-moment} now bounds the normalized \(L^p\) norm by
\(\exp(C_Dd)\), proving \eqref{p:frame-positive-bound}. The density
version and its support are inherited from that theorem. The mass-at-most-one
and absolute-continuity hypotheses have been retained explicitly.
\end{proof}

\subsection{Radial averaging and dimension-preserving propagation}

We next apply the all-radii estimate to the actual Gaussian data. A learner
may use both the row directions and their lengths and correlations. The
following factorization retains all of that radial information before
averaging it into the routing weights.

\begin{lemma}[Gaussian radial factor and frame]\label{p:radial-factor}
Let \(X\) be a \(k\)-by-\(d\) standard Gaussian matrix with \(k<d\).
Almost surely,
\[
 X=HP,\qquad H=(XX^{\mathsf T})^{1/2},\qquad P=H^{-1}X.
\]
The frame \(P\) has law \(\vartheta_{d,k}\) and is independent of \(H\).
\end{lemma}

\begin{proof}
Full row rank holds almost surely because the rank-deficient matrices are
the zero set of nonzero minors and the Gaussian law has a density. Right
multiplication of \(X\) by a fixed orthogonal matrix leaves its law
unchanged, fixes \(H\), and right-rotates \(P\). Let \(O\) be an independent
uniform orthogonal matrix. For bounded measurable \(\phi,\psi\), invariance
and averaging give
\[
 \E[\phi(H)\psi(P)]
 =\E\!\left[\phi(H)\int\psi(PO)\,d\mu(O)\right],
\]
where \(\mu\) is uniform orthogonal probability. The inner integral is
independent of \(P\): the right action on \(\mathcal V_{d,k}\) is
transitive, and \(\mu\) is invariant. Its value is
\(\int\psi\,d\vartheta_{d,k}\). This factorizes the expectation for all
bounded measurable tests, which proves the stated law and independence.
\end{proof}

\begin{proposition}[One-block all-radii propagation]\label{p:haar-block}
For sufficiently large \(d\), use \(k=q=\lfloor d/4\rfloor\) and
\(\beta=d-1-q\).
Let \(f_1,\ldots,f_N:S^{d-1}\to[0,1]\) be Borel functions satisfying
\[
 \int_{B(b,t)}f_w\,d\sigma\le Kt^\beta
 \qquad(b\in\R^d,\ t>0,\ 1\le w\le N).
\]
Let \(Q_w(X,y)\) be nonnegative Borel functions on
\(\R^{k\times d}\times\R^k\) with \(\sum_wQ_w\le1\), and define
\[
 f_{\mathrm{in}}(s)=\E_X\sum_{w=1}^N Q_w(X,Xs)f_w(s).
\]
Then
\begin{equation}\label{p:haar-block-bound}
 \int_{B(z,R)}f_{\mathrm{in}}\,d\sigma
 \le C^dN^{1/q}K R^\beta
 \qquad(z\in\R^d,\ R>0).
\end{equation}
\end{proposition}

\begin{proof}
Write \(X=HP\) as in Lemma~\ref{p:radial-factor} and average the full
radial dependence of the routing rule:
\[
 \tau_w(P,u)=\E_H Q_w(HP,Hu).
\]
These are measurable nonnegative weights with \(\sum_w\tau_w\le1\).
Independence gives the exact identity
\begin{equation}\label{p:haar-routing}
 f_{\mathrm{in}}(s)=
 \int\sum_w\tau_w(P,Ps)f_w(s)\,d\vartheta_{d,k}(P).
\end{equation}
Fix \(E=B(z,R)\), and let \(g_w(P,u)\) be the exact density in
Corollary~\ref{p:frame-n30} for \(\1_Ef_w\sigma\). After integrating
\eqref{p:haar-routing} over \(E\), pushforward integration yields
\(\int\!\int\sum_w\tau_w g_w\,du\,d\vartheta_{d,k}\).
Each \(g_w\) is supported in \(\norm{u-Pz}\le R\).
On this domain the measure with density \(\tau_w\), relative to frame
probability, label Lebesgue measure, and counting measure in \(w\), has
total mass at most \(v_kR^k\). Hölder's inequality therefore gives
\[
 \int_E f_{\mathrm{in}}\,d\sigma
 \le (v_kR^k)^{1-1/q}
 \left(\sum_w\int\!\!\int\tau_w g_w^q\,du\,d\vartheta_{d,k}\right)^{1/q}
 \le (v_kR^k)^{1-1/q}N^{1/q}\max_w\norm{g_w}_{L^q}.
\]
Substitute \eqref{p:frame-n30-bound}. The radius powers add exactly to
\(\beta\). Also \(k\ge d/8\), so \eqref{p:unit-ball-volume} gives
\[
 v_k^{1-1/q}(C\sqrt d)^{k(1-1/q)}\le C_0^d.
\]
This proves \eqref{p:haar-block-bound}. The original \(Q_w\) could use
all entries of \(H\); only after that dependence was retained did we
average it into \(\tau_w\).
\end{proof}

\subsection{Backward propagation of continuation success}

We now apply Proposition~\ref{p:haar-block} to the success probability of
a fresh continuation from each state. Its local-mass exponent survives
every block, and the resulting finite-dimensional estimate retains the
full dependence on \(M\).

\begin{proposition}[Success bound from orthonormal projections]
\label{p:haar-success-proposition}
There are absolute constants \(C\ge1\) and \(d_0\) such that the
following holds for every integer \(d\ge d_0\), every nonnegative integer
\(M\), every \(0<\epsilon\le1/10\), and every learner of deterministic
finite horizon \(T\) in Section~\ref{p:geometry-model}. Put
\[
 k=q=\lfloor d/4\rfloor,\qquad
 b=\lceil T/k\rceil,\qquad N_*=(k+2)2^M.
\]
Then
\begin{equation}\label{p:haar-success}
 \Prob\{\text{success}\}\le
 C^d\epsilon^q\bigl[C^dN_*^{1/q}\bigr]^b.
\end{equation}
The constant and dimension threshold are independent of \(M,\epsilon,T\)
and the learner's rules.
\end{proposition}

\begin{proof}
Use the fixed-seed and Borel-kernel convention of
Section~\ref{p:geometry-model}, and set \(n=d-1\) and \(\beta=n-q\).
For each terminal state and stopping index, let \(h(s)\) be the success
probability of its signal-independent output law. The cap estimates of
Lemma~\ref{p:cap-mass} give
\[
 \int_{B(z,R)}h\,d\sigma
 \le C^d\min\{R^n,\epsilon^n\}
 \le C^d\epsilon^q R^\beta
 \qquad(z\in\R^d,\ R>0).
\]
The last inequality follows from \(R\le\epsilon\) or
\(R\ge\epsilon\), respectively. Thus the all-radii coefficient for
every terminal function is at most \(K_0=C^d\epsilon^q\).

Divide the horizon into \(b\) blocks of \(k\) rows, padding the last
block with independent rows ignored by its rules. Starting at a specified
state and boundary index, define its continuation by a fresh future
execution. This definition does not condition on that state being reached.
The continuation functions are Borel and take values in \([0,1]\) by
Lemma~\ref{p:borel-kernels}.

Within a block, a destination is either a continuing state at the next
boundary or a terminal state tagged by its local stopping index. The
possible local stopping indices are \(0,1,\ldots,k\), including an
immediate stop. There are at most \((k+1)2^M\) terminal destinations and
\(2^M\) continuing destinations, hence at most \(N_*\) in total. The
global stopping index is determined by the fixed block start and this
local index, so the terminal output retains its full allowed dependence
on time.

Finite composition of the transition kernels produces Borel routing
probabilities \(Q_w(X,y)\) for these destinations, with
\(\sum_wQ_w\le1\). Their continuation functions \(f_w\) use only
future independent rows and coins; terminal destinations use the functions
\(h\) above. The incoming continuation therefore has precisely the form
\[
 f_{\mathrm{in}}(s)=\E_X\sum_wQ_w(X,Xs)f_w(s)
\]
from Proposition~\ref{p:haar-block}. The Gaussian block is fresh:
conditional on the fixed shared seed, its rows are jointly independent of
the signal and entering state, which depends only on initialization,
past pairs, and past coins.

Backward induction through the blocks multiplies the all-radii
coefficient by at most \(C^dN_*^{1/q}\) per block. This factor is at
least one, so the coefficient also bounds any earlier terminal option.
Integrate the initial continuation over \(B(0,1)\) and average over the
conditional initial-state law, which is independent of the signal and
rows. Conditional success is at most the right side of
\eqref{p:haar-success}. The bound is uniform in the fixed rules, so the
shared-seed convention gives the same bound for the original learner.
When \(T=0\), the terminal bound supplies the conclusion directly.
\end{proof}

\begin{proof}[Proof of Corollary~\ref{p:finite-memory-tradeoff} and first
proof of Theorem~\ref{p:regression}]
Put \(L=\log(1/\epsilon)\). Choose an absolute dimension threshold so
that Proposition~\ref{p:haar-success-proposition} applies and
\(k=q\ge d/8\). If success is at least \(2/3\),
\eqref{p:haar-success} implies
\[
 qL\le (b+1)d\log C
       +\frac bq\bigl(M\log2+\log(k+2)\bigr)
       +\log(3/2).
\]
Since \(b\le8T/d+1\) and \(\log(k+2)/q\) is bounded by an absolute
constant, this gives an absolute \(C_1\) such that
\begin{equation}\label{p:haar-memory-arithmetic}
 dL\le C_1(T+d)\left(1+\frac M{d^2}\right).
\end{equation}
No bound on \(M\) was used. Lemma~\ref{p:haar-linear} forces
\(T>d/8\), and hence \(T+d<9T\). Substitution into
\eqref{p:haar-memory-arithmetic} proves
Corollary~\ref{p:finite-memory-tradeoff} with an absolute constant.

Finally, if \(M(d)=o(d^2)\), then \(1+M(d)/d^2\le2\) for all
sufficiently large \(d\). The corollary proves
Theorem~\ref{p:regression}, with an absolute lower-bound constant and a
dimension threshold that may depend on the memory sequence. This proof
is uniform over the accuracy sequence and uses only within-block stopping
indices.
\end{proof}
